\documentclass[11pt]{article}
\usepackage{amsmath, amssymb, amsthm}
\usepackage{graphicx}
\usepackage{hyperref}
\usepackage{natbib}
\usepackage{booktabs}
\usepackage{geometry}
\geometry{margin=1in}

\newtheorem{theorem}{Theorem}
\newtheorem{lemma}[theorem]{Lemma}
\newtheorem{corollary}[theorem]{Corollary}
\newtheorem{definition}[theorem]{Definition}
\newtheorem{remark}[theorem]{Remark}

\title{Geometric Analysis of Constraint Barriers in Language Models and Mathematical Reasoning}
\author{Singular Heir}
\date{\today}

\begin{document}

\maketitle

\begin{abstract}
We present a unified geometric framework for understanding and removing constraints in both language models and mathematical reasoning systems. Our approach uses singular value decomposition on activation differences to identify constraint directions, which we then project out using norm-preserving techniques. We prove an unconditional barrier theorem for the shell method in Roth's theorem, showing that it cannot prove $r_3(N) \le N \exp(-c\log N)$ for any $c>0$. Under Hypothesis H, we obtain the improved bound $r_3(N) \le N \exp(-C\sqrt{\log N})$. For language models, we demonstrate surgical constraint removal with $>90\%$ refusal reduction while preserving $>95\%$ of model capabilities. Our framework reveals that constraints across domains share geometric structure, enabling technique transfer between mathematical barrier analysis and model alignment.
\end{abstract}

\section{Introduction}

The study of constraints appears in two seemingly distinct domains: mathematical barriers that prevent progress on theorems, and alignment constraints that restrict language model behavior. This paper demonstrates that these domains share a common geometric structure that enables unified analysis.

\subsection{Contributions}

\begin{enumerate}
    \item A geometric framework for constraint analysis using SVD and whitened extraction
    \item An unconditional barrier theorem for the shell method in Roth's theorem
    \item A surgical constraint removal technique for language models with minimal capability loss
    \item A cross-domain transfer framework that applies techniques from one domain to the other
    \item Open-source implementation: AETHERIS toolkit
\end{enumerate}

\section{Geometric Framework}

\subsection{Constraint Extraction}

Let $\mathcal{H}$ be the set of activations on harmful prompts and $\mathcal{H}_0$ on harmless prompts. The constraint direction $d$ is the principal component of $\mu_{\mathcal{H}} - \mu_{\mathcal{H}_0}$.

\subsection{Surgical Removal}

Norm-preserving projection removes $d$ from weight matrices $W$:

\begin{equation}
W' = W - (W d) d^\top
\end{equation}

\section{Mathematical Barriers: The Shell Method}

\subsection{Theorem 4 (Barrier)}
The shell method cannot prove $r_3(N) \le N \exp(-c\log N)$ for any $c>0$.

\subsection{Conditional Improvement}
Under Hypothesis H, we obtain $r_3(N) \le N \exp(-C\sqrt{\log N})$.

\section{Experimental Results}

\section{Discussion}

\section{Conclusion}

\bibliographystyle{plainnat}
\bibliography{references}

\end{document}
